% Part: first-order-logic
% Chapter: syntax-and-semantics
% Section: main-operator

\documentclass[../../../include/open-logic-section]{subfiles}

\begin{document}

\olfileid{fol}{syn}{mai}

\olsection{د فارمول \printtoken{S}{main operator}}

\begin{explain}
ډېر ځله د هغه وروستي عملګر په باب خبرې کول ګټور وي چې د
!!a{formula}~$!A$ په جوړولو کښې کارول شوے وي۔ دا عملګر د~$!A$
\emph{اصلي عملګر} بلل کېږي۔ په شهودي ډول دا د~$!A$ ``تر ټولو دباندې''
عملګر دے۔ مثلاً د $\lnot !A$ اصلي عملګر $\lnot$ دے، د $(!A \lor
!B)$ اصلي عملګر $\lor$ دے، او همداسې نور۔
\end{explain}


\begin{defn}[!!^{main operator}]
\ollabel{def:main-op}
د !!a{formula}~$!A$ \emph{!!{main operator}} په لاندې ډول تعريفېږي:
\begin{enumerate}
\item \indcase*{!A}{!A}{$\indfrm$ هېڅ !!{main operator} نۀ لري۔}

\tagitem{prvNot}{\indcase{!A}{\lnot !B}{د $\indfrm$ !!{main operator}
  سمبول~$\lnot$ دے۔}}{}

\tagitem{prvAnd}{\indcase{!A}{(!B \land !C)}{د $\indfrm$
  !!{main operator} سمبول~$\land$ دے۔}}{}

\tagitem{prvOr}{\indcase{!A}{(!B \lor !C)}{د $\indfrm$
  !!{main operator} سمبول~$\lor$ دے۔}}{}

\tagitem{prvIf}{\indcase{!A}{(!B \lif !C)}{د $\indfrm$
  !!{main operator} سمبول~$\lif$ دے۔}}{}

\tagitem{prvIff}{\indcase{!A}{(!B \liff !C)}{د $\indfrm$
  !!{main operator} سمبول~$\liff$ دے۔}}{}

\tagitem{prvAll}{\indcase{!A}{\lforall[x][!B]}{د $\indfrm$
  !!{main operator} سمبول~$\lforall$ دے۔}}{}

\tagitem{prvEx}{\indcase{!A}{\lexists[x][!B]}{د $\indfrm$
  !!{main operator} سمبول~$\lexists$ دے۔}}{}
\end{enumerate}
\end{defn}

په هر حالت کښې زموږ موخه په فارمول کښې د !!{main operator} همدا ښودل
شوے ځانګړے \emph{مورد} دے۔ مثلاً فارمول $((!D \lif !E) \lif (!E
\lif !D))$ د $(!B \lif !C)$ بڼه لري، چې $!B$ پکښې $(!D \lif !E)$
او $!C$ پکښې $(!E \lif !D)$ دے؛ نو د $\lif$ دويم مورد ئې
!!{main operator} دے۔

\begin{explain}
دا د هغې تابع يو \emph{بازګشتي} تعريف دے چې غير اټومي
!!{formula}s د هغو د !!{main operator} موردونو ته لېږي۔ څرنګه چې
!!{formula}s په استقرايي ډول تعريف شوي دي، هر !!{formula}~$!A$ په
\olref{def:main-op} کښې له ورکړل شوو حالتونو څخه يو پوره کوي۔ دا
تضمينوي چې د هر غير اټومي !!{formula}~$!A$ دپاره !!a{main
  operator} موجود دے۔ څرنګه چې هر !!{formula} له دغو شرطونو څخه
يوازې يو پوره کوي، او په هر حالت کښې هغه کوچني !!{formula}s په
يوازيني ډول ټاکل شوي دي چې $!A$ ترې جوړ شوے، د~$!A$ د !!{main
  operator} مورد هم يوازينے دے؛ نو په رښتيا مو يوه تابع تعريف کړې ده۔
\end{explain}

موږ !!{formula}s په
\olref{tab:main-op} کښې د هغو د !!{main operator} د سمبول له مخې په
ورکړل شوو نومونو يادوو۔
\iftag{defNot,defOr,defAnd,defIf,defIff,defTrue,defFalse,defEx,defAll}
{ خو په ياد ولرئ چې تعريف شوي عملګرونه په رسمي ډول په !!{formula}s
کښې نۀ راځي۔ هغوئ يوازې لنډونونه دي، نو په رسمي ډول د يوه فارمول
اصلي عملګر نۀ شي کېدے۔ ځکه د ټولو !!{formula}s په باب ثبوتونو کښې
جلا چلند ورسره پکار نۀ دے۔}

\begin{table}[!h]
\centering
\begin{tabular}{c | c | c}
!!^{main operator} & د !!{formula} ډول & بېلګه\\
\hline
هېڅ & اټومي (!!{formula}) &
\iftag{prvFalse}{$\lfalse$,}{}
\iftag{prvTrue}{$\ltrue$,}{}
$\Atom{R}{t_1, \dots, t_n}$,
$\eq[t_1][t_2]$\\
$\lnot$ & نفي & $\lnot !A$ \\
$\land$ & عطف & $(!A \land !B)$ \\
$\lor$ & فصل & $(!A \lor !B)$ \\
$\lif$ & !!{conditional} & $(!A \lif !B$) \\
$\liff$ & !!{biconditional} & $(!A \liff !B)$ \\
$\lforall[][]$ & عمومي (!!{formula})& $\lforall[x][!A]$ \\
$\lexists[][]$ & وجودي (!!{formula})& $\lexists[x][!A]$
\end{tabular}
\caption{اصلي عملګر او د !!{formula}s نومونه}
\ollabel{tab:main-op}
\end{table}
\textbf{د ژباړې د نسخې سمون (OLFOL-007):} سرچينه د عطف او فصل د
بېلګو تړونکي قوسونه د رياضي له چاپېريال څخه بهر ږدي؛ دلته دواړه قوسونه
له خپلو فارمولونو سره يو ځاے شوي دي۔

\end{document}
